\section{Turbulence closures}
\label{sec:turb}

Turbulence is the case with no small parameter: eddies of all sizes between
the mean-flow scale and the dissipation scale are present, and none is
separated from the scale being resolved. The closures here are
\emph{models}. They are included because the supernova profile needs them,
because most of the vocabulary for testing closures (a priori and a
posteriori tests, realizability) was developed for them, and because the
contrast with Section~\ref{sec:asym} is instructive. The standard reference
is Pope \cite{Pope}.

\subsection{Eddy viscosity}

Postulate the form of \eqref{eq:NSF} and model the coefficient. For the
Reynolds stress (Boussinesq, 1877),
\begin{equation}
  \avg{u_i'u_j'}-\tfrac23k\,\delta_{ij}=-2\nu_t\,\bar S_{ij},
  \qquad
  \bar S_{ij}=\tfrac12(\partial_i\bar u_j+\partial_j\bar u_i),
  \qquad k=\tfrac12\avg{u_i'u_i'} .
  \label{eq:boussinesq}
\end{equation}
The assumption is tensorial: the anisotropic stress is aligned with the
local mean strain. The eddy viscosity $\nu_t$ is then supplied at one of
several levels.
\begin{itemize}
  \item \textbf{Mixing length} (Prandtl, 1925):
    $\nu_t=\ell_m^2\,|\partial\bar u/\partial y|$ with a prescribed length
    $\ell_m$. This is the kinetic estimate $\nu\sim\vt\lambda$ with eddies
    in place of molecules. The analogy is not a derivation, because the
    ``mean free path'' of an eddy is comparable to the scale of the mean
    flow.
  \item \textbf{Two-equation models.} Transport two turbulence quantities
    and form $\nu_t$ from them. In the $k$--$\varepsilon$ model
    \cite{LaunderSpalding74},
    \begin{align}
      \nu_t &= C_\mu\,\frac{k^2}{\varepsilon}, \nonumber\\
      \frac{\bar Dk}{\bar Dt} &= P-\varepsilon
         +\nabla\cdot\Bigl[\Bigl(\nu+\frac{\nu_t}{\sigma_k}\Bigr)\nabla k\Bigr],
         \qquad P=2\nu_t\,\bar S_{ij}\bar S_{ij}, \\
      \frac{\bar D\varepsilon}{\bar Dt} &= C_{\varepsilon1}\frac{\varepsilon}{k}P
         -C_{\varepsilon2}\frac{\varepsilon^2}{k}
         +\nabla\cdot\Bigl[\Bigl(\nu+\frac{\nu_t}{\sigma_\varepsilon}\Bigr)
           \nabla\varepsilon\Bigr], \nonumber
    \end{align}
    with standard constants $C_\mu=0.09$, $C_{\varepsilon1}=1.44$,
    $C_{\varepsilon2}=1.92$, $\sigma_k=1.0$, $\sigma_\varepsilon=1.3$.
\end{itemize}
The constants are calibrated on canonical flows: $C_{\varepsilon2}$ from
the decay of grid turbulence, $C_\mu$ from the ratio of shear stress to $k$
in the logarithmic layer, and the others from requiring the logarithmic
law. The extra equations do not remove the closure problem. The $k$
equation is derived and then modelled term by term, and the $\varepsilon$
equation is largely constructed by analogy.

\emph{Realizability.} Equation \eqref{eq:boussinesq} gives
$\avg{u_1'^2}=\tfrac23k-2\nu_t\bar S_{11}$, which is negative when
$\bar S_{11}>k/(3\nu_t)$. A true Reynolds stress is a covariance and is
positive semidefinite. Realizable variants limit $C_\mu$ in strong strain.

\subsection{Second-moment closures}

Reynolds-stress transport models evolve all six components of $R_{ij}$,
\[
  \frac{\bar DR_{ij}}{\bar Dt}=P_{ij}+\Phi_{ij}-\varepsilon_{ij}+D_{ij} .
\]
Production $P_{ij}$ is closed. The pressure--strain redistribution
$\Phi_{ij}$, the dissipation tensor and the turbulent transport $D_{ij}$
are not. The closure has moved one level up the hierarchy, exactly as with
Grad's equations relative to Navier--Stokes--Fourier. Algebraic stress
models sit between the two levels: they expand the anisotropy in a tensor
basis built from mean strain and rotation, with scalar coefficients that
depend on invariants. That representation is the one tensor-basis neural
networks later adopted (Section~\ref{sec:learned}).

\subsection{Large-eddy simulation}

LES resolves the large eddies and models the subgrid stress \eqref{eq:sgs}.
Two families of model reflect two goals.

\emph{Functional models} aim at the right net energy transfer to the
unresolved scales. The Smagorinsky model \cite{Smagorinsky63} is
\begin{equation}
  \tau_{ij}-\tfrac13\tau_{kk}\delta_{ij}=-2\nu_t\bar S_{ij},
  \qquad
  \nu_t=(C_s\Delta)^2|\bar S|,
  \qquad |\bar S|=\sqrt{2\bar S_{ij}\bar S_{ij}} .
  \label{eq:smag}
\end{equation}
Its transfer $-\tau_{ij}\bar S_{ij}=\nu_t|\bar S|^2\ge0$ always removes
energy from the resolved field. Requiring that it remove energy at the
rate of the Kolmogorov cascade gives
$C_s=\pi^{-1}(2/3C_K)^{3/4}\approx0.17$ for $C_K=1.5$
(Exercise~\ref{ex:smag}). In shear flows smaller values are needed, which
is the motivation for the dynamic procedure.

\emph{The dynamic procedure} \cite{Germano91,Lilly92} computes the
coefficient from the resolved field. Apply a second, wider test filter
(hat) and let $T_{ij}$ be the subgrid stress at the test level. Then
\begin{equation}
  L_{ij}\equiv T_{ij}-\widehat{\tau_{ij}}
  =\widehat{\bar u_i\bar u_j}-\hat{\bar u}_i\hat{\bar u}_j
\end{equation}
is computable from the resolved velocity (the Germano identity).
Inserting \eqref{eq:smag} at both levels with the same coefficient
$C=C_s^2$ gives $L_{ij}^{d}=C\,M_{ij}$ with
$M_{ij}=2\Delta^2\widehat{|\bar S|\bar S_{ij}}
-2\hat\Delta^2|\hat{\bar S}|\hat{\bar S}_{ij}$, and the least-squares
solution is $C=\avg{L_{ij}M_{ij}}/\avg{M_{ij}M_{ij}}$. The idea is general:
the closure is required to be consistent across two levels of
coarse-graining, and its free parameter is fixed by that requirement.

\emph{Structural models} aim at the stress tensor itself. The gradient
model $\tau_{ij}\approx\tfrac1{12}\Delta^2\,\partial_k\bar u_i\,\partial_k\bar u_j$
is the leading term of a Taylor expansion of \eqref{eq:sgs}; scale
similarity and approximate deconvolution are related. In \emph{implicit}
LES no stress is written and the dissipation of the numerical scheme is the
closure, so the discretization is part of the model.

\subsection{The lesson of a priori tests}

Clark, Ferziger and Reynolds \cite{Clark79} filtered direct simulations and
compared modelled and exact subgrid stresses point by point. The
Smagorinsky stress correlates poorly with the exact one, and structural
models correlate well. Embedded in a simulation the ranking reverses: the
Smagorinsky model gives stable runs with reasonable statistics, and purely
structural models tend to be unstable because they do not dissipate enough.
The embedded behaviour depends on the net energy transfer, which a
pointwise correlation does not measure. This is the origin of the rule that
a closure must be tested both ways, and it is inequality \eqref{eq:gronwall}
read in the other direction.

\subsection{Statistical closure theories}
\label{sec:statclosure}

For homogeneous turbulence one can write the hierarchy for the spectrum
$E(k,t)$ directly: its equation involves triple correlations, whose
equation involves fourth-order ones.
\begin{itemize}
  \item \textbf{Quasi-normal.} Discard the fourth cumulant, as in
    Example~\ref{ex:langevin}. The resulting equations develop negative
    energy spectra: the closure is not realizable.
  \item \textbf{EDQNM} \cite{Orszag70}. Add a damping of the triple
    correlations at an eddy-turnover rate and make the equation Markovian.
    Realizable, with a free constant.
  \item \textbf{Direct-interaction approximation}
    \cite{Kraichnan59}. A self-consistent closure for the two-time
    correlation and the response function, with no free constant. It is the
    exact statistics of a model system and therefore realizable. In its
    Eulerian form it does not reproduce the Kolmogorov spectrum, because it
    does not separate sweeping of small eddies by large ones from their
    distortion; Lagrangian versions repair this.
\end{itemize}
The response function introduced by the direct-interaction approximation is
the same object as the memory kernel of Section~\ref{sec:memory}, and the
plasma versions (resonance broadening, Section~\ref{sec:waves}) are
closures of this type.
